% !PRINTED 406
\weberpage{406}
%<<CONT
in which $m$ runs through all the positive integers relatively prime to $Q$,
and $n$, for each $m$, runs through the divisors of $m$.

If we arrange the right side of (2) according to $n$, then for a
given $n$ all the multiples of $n$ occur among the $m$, and we
can also set:

\weq{3}{\Sigma F(N\mathfrak{a}) = \overset{m}{\Sigma}\overset{n}{\Sigma}(\Delta,n)\,F(mn),}

in which now $m$ and $n$, independently of one another, assume all
natural numerical values that have no divisor in common with $Q$.

If in particular $F$ is such that

\weq{4}{F(mn) = F(m)\,F(n)}

holds, then we can also represent (3) in this way:

\weq{5}{\overset{\mathfrak{a}}{\Sigma}F(N\mathfrak{a}) = \overset{m}{\Sigma}F(m)\overset{n}{\Sigma}(\Delta,n)\,F(n).}

In the sums over $m$ and over $n$ those numbers are excluded
that have a divisor in common with $Q$.

Likewise, on the left side those ideals are excluded
that are not relatively prime to $Q$.

We can apply formula (5) first of all in order to determine the
\gesperrt{class number} from the formulae of the preceding section.
In this way one also obtains the ratio of the class number
in the \gesperrt{field} to the class number with respect to the order $[Q]$. But since
we have already determined this ratio in another way
(§ 100), we shall here restrict ourselves to the \gesperrt{class number of the
 field}, i.\,e. we shall set $Q = 1$.

If in (5) we set $F(z) = z^{-s}$, then the presupposed
unconditional convergence takes place as long as $s > 1$. But if we
multiply by $s - 1$ and let $s$ pass over into $1$, we can
make use of § 111, (12).

We apply this formula in the first version, in
which $\mathfrak{m}$ runs through the ideals of the class $M$ coprime to $Q$,
and $r$ denotes the prime factors of $Q$. These therefore drop out
entirely for $Q = 1$, and $\mathfrak{m}$ in the summation formula:

\weq{6}{\begin{aligned}
\sum\frac{s-1}{(N\mathfrak{m})^s} &= \frac{1}{\sqrt{\Delta}}\log\varepsilon, & \Delta&>0\\
&= \frac{\pi}{\sqrt{-\Delta}}, & \Delta&<0
\end{aligned}}

is no longer subject to any restriction with regard to divisibility
.
